% ARCHIVED at v3.88 (2026-09-24) from chapters/06_results.tex -- author: "6.4 Comparison across Environments at least could show all the things besides the scurve now". The v3.60 lock text and the pre-reconciliation section (corridor rows from the withdrawn v2 corpus, no UAV-pillars rows).
% Kept verbatim for rollback; not part of the build.

% =============================================================================
% =============================================================================
% 🔒 DO NOT EDIT (author, 2026-09-21 / v3.60)
%
% This section summarises 6.1, 6.2 and 6.3. It is the LAST thing to be written,
% because every number in it is a restatement of a number from those three, and
% a summary written over sections that are still moving has to be rewritten each
% time one of them moves.
%
% Leave it exactly as it stands until 6.1, 6.2 and 6.3 have each been verified
% against the corpus and locked by the author. Then reconcile this section with
% them in one pass. Do not "fix" a cell here on its own -- if a number here
% disagrees with its source section, the source section is what to check.
% =============================================================================
\section{Comparison across Environments}\label{sec:res:summary}
% =============================================================================

This section recaps the preceding ones and adds no new result.

\subsection{Generative Models}\label{sec:res:summary:models}

\autoref{tab:summary-models} names, per environment, the generative model that comes out best and the
measure that decides it. Analytic average-velocity matching is best, or shares first place, in each of the
three environments that yield a result.

\begin{table}[htpb]
  \caption[Generative models across environments]{The best generative model per environment, as found in
  Sections~\ref{sec:res:state}--\ref{sec:res:uav}. Temporal U-Net, 3.96--4.0\,M parameters, throughout.}
  \label{tab:summary-models}
  \centering
  \small
  \begin{tabularx}{\linewidth}{@{}p{2.3cm} L p{2.6cm} L@{}}
    \toprule
    Environment & Deciding measure & Best & The others \\
    \midrule
    D3IL-avoiding & control steps and time at equal success & MeanFM, CI-MeanFM & FM, then Diffusion \\
    D3IL-aligning & final distance of the box & MeanFM & CI-MeanFM, FM and Diffusion level \\
    UAV-corridor  & success on a collision-free flight & MeanFM, CI-MeanFM ($12/12$) & FM succeeds but not collision-free; Diffusion collision-free but never succeeds \\
    UAV-s-curve   & --- & none & limited by the tracking controller \\
    \bottomrule
  \end{tabularx}
\end{table}

\subsection{Projection Methods}\label{sec:res:constraints}

\autoref{tab:summary-projection} names the projection method each environment selects. Per-step
projection is selected where the budget is too small for endpoint projection to act and on UAV-corridor;
endpoint projection is selected on D3IL-aligning, the task run at a budget that gives it
guiding steps.

\begin{table}[htpb]
  \caption[Projection methods across environments]{The selected projection method per environment, as
  found in Sections~\ref{sec:res:state}--\ref{sec:res:uav}.}
  \label{tab:summary-projection}
  \centering
  \small
  \begin{tabularx}{\linewidth}{@{}p{2.3cm} r p{2.3cm} L@{}}
    \toprule
    Environment & $\nfe$ & Selected & Reason \\
    \midrule
    D3IL-avoiding & 1  & per-step & no guiding step for endpoint projection \\
    D3IL-aligning & 20 & endpoint & all ten contexts free of violations \\
    UAV-corridor  & 3  & per-step & endpoint projection succeeds collision-free on at most 7 of 12 \\
    UAV-s-curve   & ---& none     & no flight succeeds within the constraints \\
    \bottomrule
  \end{tabularx}
\end{table}

\subsection{Conclusion}\label{sec:res:summary:conclusion}

% v3.68d (author, 2026-09-23): this paragraph is the one addition made under the 🔒 lock, at the
% author's explicit request for a sub-conclusion beside the table.
Read across the environments, the generative models order the same way where the task separates them at
all. On D3IL-avoiding they do not separate on the outcome --- CI-MeanFM $\approx$ MeanFM $\approx$ FM
$\approx$ diffusion --- and separate on the price by a factor of thirty in favour of the flow-based models.
On D3IL-aligning they separate on the outcome, MeanFM $\gg$ CI-MeanFM $>$ FM $\approx$ diffusion. On the
quadrotor scenes the reading is pending the runs of \autoref{sec:res:uav:conclusion}. The projector orders
the other way round by task: per-step projection at a small budget where the task is solved at one or two
evaluations, endpoint projection at the large budget the alignment task needs.
\autoref{tab:summary-combinations} collects the selected combination of each environment.

\begin{table}[htpb]
  \caption[Selected model--projection combinations]{The selected combination per environment.}
  \label{tab:summary-combinations}
  \centering
  \small
  \begin{tabularx}{\linewidth}{@{}p{2.3cm} L L@{}}
    \toprule
    Environment & Combination & Result \\
    \midrule
    D3IL-avoiding
      & MeanFM or CI-MeanFM, $\nfe=1$, per-step projection, temporal consistency
      & S\&C $0.967$ and $1.000$ at $18.3$ and $18.1$\,ms per step; five seeds, the protocol of \ac{DPCC} \\
    D3IL-aligning
      & MeanFM, $\nfe=20$, endpoint projection, random selection
      & $10/10$ contexts free of violations, box moved in $8/10$ \\
    UAV-corridor
      & MeanFM, $\nfe=3$, per-step projection, temporal consistency
      & success, collision-free, in $12/12$ flights \\
    UAV-s-curve
      & none
      & limited by the tracking controller \\
    \bottomrule
  \end{tabularx}
\end{table}

